% !TEX root = ../matan.tex

\section{Интегрирование и дифференцирование степенных рядов. Аналитические функции.}

\begin{Theorem}{Почленное интегрирование и дифференцирование}{}
	Пусть ряд $\sum_{n=0}^{\infty} a_n (x - x_0)^n$ имеет радиус сходимости $R > 0$. Тогда внутри интервала сходимости $(x_0 - R,\, x_0 + R)$ его сумма $S(x) = \sum_{n=0}^{\infty} a_n (x - x_0)^n$ дифференцируема, и ряд можно почленно дифференцировать и интегрировать (на любом $[x_0 - r,\, x_0 + r]$ с $0 < r < R$):
	\[
	S'(x) = \sum_{n=1}^{\infty} n\,a_n (x - x_0)^{n-1}, \qquad
	\int_{x_0}^{x} S(t)\,dt = \sum_{n=0}^{\infty} \frac{a_n}{n + 1} (x - x_0)^{n+1}.
	\]
	Оба полученных ряда имеют тот же радиус сходимости $R$.
\end{Theorem}

\begin{proof}
	Без ограничения общности $x_0 = 0$. Продифференцированный ряд есть $\sum_{n \ge 1} n\,a_n x^{n-1} = \sum_{m \ge 0} (m + 1) a_{m+1} x^m$. По формуле Коши--Адамара его радиус
	\[
	R_1 = \frac{1}{\limsup_{m \to \infty} |(m + 1) a_{m+1}|^{1/m}}.
	\]
	Поскольку $(m + 1)^{1/m} \to 1$, верхний предел не меняется: $\limsup_m |(m+1)a_{m+1}|^{1/m} = \limsup_m |a_{m+1}|^{1/m}$. Далее, $|a_{m+1}|^{1/m} = \bigl(|a_{m+1}|^{1/(m+1)}\bigr)^{(m+1)/m}$, а показатель $(m+1)/m \to 1$, поэтому
	\[
	\limsup_{m \to \infty} |a_{m+1}|^{1/m} = \limsup_{m \to \infty} |a_{m+1}|^{1/(m+1)} = \limsup_{n \to \infty} |a_n|^{1/n} = \frac{1}{R}.
	\]
	Значит, $R_1 = R$. Аналогично проинтегрированный ряд $\sum a_n x^{n+1}/(n+1)$ имеет радиус $R$, так как деление на $n + 1$ не влияет на $\limsup |a_n/(n+1)|^{1/n} = \limsup|a_n|^{1/n}$.

	Зафиксируем $0 < r < R$. На $[-r, r]$ как исходный ряд, так и ряд производных сходятся равномерно (степенной ряд сходится равномерно на любом замкнутом круге строго меньшего радиуса). Все частичные суммы~--- многочлены, то есть непрерывно дифференцируемы. По теореме о почленном дифференцировании равномерно сходящегося ряда сумма $S$ дифференцируема на $[-r, r]$ и $S'(x) = \sum_{n \ge 1} n\,a_n x^{n-1}$. Так как $r < R$ произвольно, это верно на всём $(-R, R)$. Формула для интеграла получается почленным интегрированием равномерно сходящегося на $[-r, r]$ ряда. % опирается на вопросы 28-29, 31
\end{proof}

\begin{Corollary}{Бесконечная дифференцируемость}{}
	Сумма степенного ряда бесконечно дифференцируема внутри интервала сходимости, причём коэффициенты восстанавливаются по производным в центре:
	\[
	a_k = \frac{S^{(k)}(x_0)}{k!}.
	\]
\end{Corollary}

\begin{proof}
	Применяя предыдущую теорему $k$ раз (каждое дифференцирование сохраняет радиус $R$), получаем
	\[
	S^{(k)}(x) = \sum_{n \ge k} a_n\, n(n-1)\cdots(n - k + 1)\,(x - x_0)^{n-k}.
	\]
	При $x = x_0$ все слагаемые с $n > k$ обращаются в нуль, остаётся член $n = k$: $S^{(k)}(x_0) = a_k\, k!$, откуда $a_k = S^{(k)}(x_0)/k!$. Возможность дифференцировать сколько угодно раз и означает бесконечную дифференцируемость.
\end{proof}

\begin{Definition}{Аналитическая функция}{}
	Функция $f$, определённая в окрестности точки $x_0$, называется \textbf{аналитической} в $x_0$, если в некоторой окрестности этой точки она представима сходящимся степенным рядом
	\[
	f(x) = \sum_{n=0}^{\infty} a_n (x - x_0)^n.
	\]
	Функция аналитична на открытом множестве $U$, если она аналитична в каждой его точке. % [восстановлено] определение аналитичности в ticket07 явно не дано
\end{Definition}

\begin{Remark}{}{}
	По предыдущему следствию всякая аналитическая функция бесконечно дифференцируема, а сумма степенного ряда аналитична во всём своём круге сходимости. Обратное включение неверно (см. вопрос~33): класс $C^\infty$ строго шире класса аналитических функций.
\end{Remark}
